\documentclass[11pt]{article}
\usepackage[T1]{fontenc}
\usepackage[margin=1in]{geometry}
\usepackage{float}
\usepackage{caption}
\usepackage{hyperref}
\hypersetup{colorlinks=true,linkcolor=blue}

\floatstyle{ruled}
\newfloat{protocol}{H}{lop}
\floatname{protocol}{Protocol}

\title{Bench Protocols for the Teaching Laboratory}
\author{Teaching Laboratory}
\date{}

\begin{document}
\maketitle
\tableofcontents
\listoffigures
\listoftables
\listof{protocol}{List of Protocols}

\section{Introduction}
\label{sec:intro}

This manual collects three bench protocols used in the teaching laboratory.
Each protocol is typeset as a custom \texttt{protocol} float placed with the
\textsf{float} package's \texttt{[H]} specifier, so it appears exactly where
it is written -- see Protocol~\ref{pr:calibrate} and
Protocol~\ref{pr:titrate} -- while ordinary figures and tables still float.
The calibration curve of Figure~\ref{fig:curve} and the reagent stock of
Table~\ref{tab:reagents} are cross-referenced forward from this introduction,
so all of these numbers are unresolved on the first pass and settle together
with the three lists (figures, tables, protocols) on later passes.

\begin{figure}[H]
  \centering
  \rule{0.7\textwidth}{3cm}
  \caption{Calibration curve of the bench spectrometer}
  \label{fig:curve}
\end{figure}

\begin{protocol}[H]
  \caption{Calibrating the spectrometer}
  \label{pr:calibrate}
  Warm up the lamp for ten minutes. Record the dark spectrum, then the
  spectrum of the reference cell. Fit the affine correction of
  Section~\ref{sec:notes}; the fitted curve should match
  Figure~\ref{fig:curve}.
\end{protocol}

\section{Reagents}
\label{sec:reagents}

Table~\ref{tab:reagents} lists the stock solutions prepared for the term.
The titration of Protocol~\ref{pr:titrate} consumes the potassium
permanganate entry, and its endpoint color is compared against the reference
swatches of Figure~\ref{fig:swatch}.

\begin{table}[H]
  \centering
  \begin{tabular}{|l|c|c|}
    \hline
    Reagent & Concentration & Volume \\
    \hline
    KMnO$_4$ & 0.02\,mol/L & 2\,L \\
    Oxalic acid & 0.05\,mol/L & 1\,L \\
    Sulfuric acid & 1\,mol/L & 500\,mL \\
    \hline
  \end{tabular}
  \caption{Stock solutions for the term}
  \label{tab:reagents}
\end{table}

\begin{protocol}[H]
  \caption{Standardizing the permanganate}
  \label{pr:titrate}
  Pipette 25\,mL of oxalic acid into a flask, add 10\,mL of sulfuric acid,
  warm to $60^{\circ}$C, and titrate with KMnO$_4$ from
  Table~\ref{tab:reagents} to the first permanent pink.
\end{protocol}

\begin{figure}[H]
  \centering
  \rule{0.6\textwidth}{2.5cm}
  \caption{Reference endpoint swatches for the titration}
  \label{fig:swatch}
\end{figure}

\section{Notes}
\label{sec:notes}

The affine correction used in Protocol~\ref{pr:calibrate} is
$c = a\,r + b$ with the coefficients posted on the laboratory wall chart.
Figure~\ref{fig:curve}, Figure~\ref{fig:swatch}, Table~\ref{tab:reagents},
Protocol~\ref{pr:calibrate}, and Protocol~\ref{pr:titrate} are all cited
before they appear, so the lists of figures, tables, and protocols grow
between passes and float placement shifts accordingly.

\begin{protocol}[H]
  \caption{Shutting down the bench}
  \label{pr:shutdown}
  Rinse glassware, return stock bottles to Table~\ref{tab:reagents} storage,
  switch off the lamp of Figure~\ref{fig:curve}, and sign the logbook.
\end{protocol}

\end{document}
